% part: intuitionistic-logic
% chapter: semantics
% section: topological-semantics

\documentclass[../../../include/open-logic-section]{subfiles}

\begin{document}

\olfileid{int}{sem}{top}

\olsection{الدلالات الطوبولوجية}

من الطرق الأخرى لتقديم دلالات للمنطق الحدسي استعمال المفهوم الرياضي
للطوبولوجيا.

\begin{defn}
  لتكن~$X$ مجموعة. \emph{الطوبولوجيا على~$X$} مجموعة
  $\Top{O}\subseteq\Pow{X}$ تستوفي الخواص الآتية. وتسمى
  !!{element}s~$\Top{O}$ \emph{المجموعات المفتوحة} في الطوبولوجيا.
  وتسمى المجموعة~$X$ مع~$\Top{O}$ \emph{فضاءً طوبولوجيًا}.
  \begin{enumerate}
  \item المجموعة الخالية والفضاء كله مفتوحان:
    $\emptyset$، $X \in \Top{O}$.
  \item المجموعات المفتوحة مغلقة تحت التقاطعات المنتهية: إذا كان
    $U$، $V \in \Top{O}$، كان~$U \cap V \in \Top{O}$.
  \item المجموعات المفتوحة مغلقة تحت الاتحادات الاعتباطية: إذا كان
    $U_i \in \Top{O}$ لكل~$i \in I$، كان
    $\bigcup \Setabs{U_i}{i \in I} \in \Top{O}$.
  \end{enumerate}
\end{defn}

قد نكتب~$X$ اختصارًا للفضاء الطوبولوجي إذا أمكن استنتاج مجموعة المجموعات
المفتوحة من السياق. لكن ينبغي التنبه إلى أن~$X$ نفسها لا تسمى طوبولوجيا:
فالطوبولوجيا هي مجموعة المجموعات المفتوحة، أما~$X$ بعد تزويدها بتلك
المجموعات فهي فضاء طوبولوجي.

\begin{defn}
  \emph{النموذج الطوبولوجي} للمنطق القضي الحدسي ثلاثي
  $\mModel{X}=\tuple{X,\Top{O},V}$، حيث~$\Top{O}$ طوبولوجيا على~$X$،
  و$V$ دالة تسند إلى كل متغير قضي مجموعة مفتوحة في~$\Top{O}$.

  وبالنظر إلى نموذج طوبولوجي~$\mModel{X}$، يمكننا تعريف~$\Prop{X}{!A}$
  استقرائيًا كما يأتي:
  \begin{enumerate}
  \item $\Prop{X}{\lfalse}=\emptyset$؛
  \item $\Prop{X}{p}=V(p)$؛
  \item $\Prop{X}{!A \land !B}=\Prop{X}{!A}\cap\Prop{X}{!B}$؛
  \item $\Prop{X}{!A \lor !B}=\Prop{X}{!A}\cup\Prop{X}{!B}$؛
  \item $\Prop{X}{!A \lif !B}=\Interior{(X\setminus\Prop{X}{!A})\cup
    \Prop{X}{!B}}$.
  \end{enumerate}
  وهنا~$\Interior{V}$ هي الدالة التي ترسل مجموعة~$V \subseteq X$ إلى
  \emph{داخلها}، أي إلى اتحاد جميع المجموعات المفتوحة التي تحتويها.
  وبعبارة أخرى:
  \[
  \Interior{V}=\bigcup\Setabs{U}{U \subseteq V \text{ و } U \in \Top{O}}.
  \]
\end{defn}

لاحظ أن داخل أي مجموعة مفتوح دائمًا، لأنه اتحاد مجموعات مفتوحة. ومن ثم تكون
$\Prop{X}{!A}$ دائمًا مجموعة مفتوحة.

مع أن الدلالات الطوبولوجية شديدة التجريد، توجد طرائق للتفكير فيها قد توضح
الدافع إليها. افترض أن~!!{element}s، أو «نقاط»،~$X$ نقاط يمكن تقويم
العبارات عندها. فمجموعة جميع النقاط التي تصدق عندها~$!A$ هي القضية التي
تعبر عنها~$!A$. وليس كل مجموعة من النقاط قضية محتملة؛ بل وحدها
!!{element}s~$\Top{O}$ كذلك. ويكون~$!A \Entails !B$ إذا وفقط إذا كانت
$!B$ صادقة عند كل نقطة تصدق عندها~$!A$؛ أي إذا كان
$\Prop{X}{!A}\subseteq\Prop{X}{!B}$ لكل~$X$. ولا تصدق العبارة
المحال~$\lfalse$ أبدًا، ولذلك~$\Prop{X}{\lfalse}=\emptyset$.

كيف يجب أن ترتبط القضايا التي تعبر عنها~$!B \land !C$ و$!B \lor !C$
و$!B \lif !C$ بالقضايا التي تعبر عنها~$!B$ و$!C$ لكي تصح القوانين
الحدسية؛ أي لكي يكون~$!A \Proves !B$ إذا وفقط إذا كان
$\Prop{X}{!A}\subseteq\Prop{X}{!B}$؟ نشترط~$\lfalse \Proves !A$ لكل
$!A$، وهو متحقق لأن~$\emptyset\subseteq U$ لكل~$U$. ولأن
$!B \land !C \Proves !B$، نشترط
$\Prop{X}{!B \land !C}\subseteq\Prop{X}{!B}$، وبالمثل
$\Prop{X}{!B \land !C}\subseteq\Prop{X}{!C}$. وأكبر مجموعة تستوفي
$W \subseteq U$ و$W \subseteq V$ هي~$U \cap V$. وبالعكس،
$!B \Proves !B \lor !C$ و$!C \Proves !B \lor !C$، ولذلك نشترط
$\Prop{X}{!B}\subseteq\Prop{X}{!B \lor !C}$ و
$\Prop{X}{!C}\subseteq\Prop{X}{!B \lor !C}$. وأصغر مجموعة~$W$ بحيث
$U \subseteq W$ و$V \subseteq W$ هي~$U \cup V$.

تعريف~$\lif$ دقيق: فـ~$!A \lif !B$ تعبر عن أضعف قضية تستلزم~$!B$ حين
تُجمع مع~$!A$. ويتضح أن~$!A \lif !B$ مع~$!A$ تستلزم~$!B$ من
$(!A \lif !B)\land !A \Proves !B$. ولذلك ينبغي أن تكون
$\Prop{X}{!A \lif !B}$ أكبر مجموعة مفتوحة بحيث
$\Prop{X}{!A \lif !B}\cap\Prop{X}{!A}\subseteq\Prop{X}{!B}$، وهذا
يفضي إلى تعريفنا.

\end{document}
